% ------------------------------------------------------------------------
% RELATÓRIO TÉCNICO OPENCODE
% Formato: ABNT NBR 6022:2018 (Artigo em Publicação Técnico-Científica)
% Compilação: lualatex → biber → lualatex → lualatex
% ------------------------------------------------------------------------
\documentclass[
  brazil,
  sumario=tradicional,
  10pt
]{abntex2}

% --- Pacotes fundamentais ---
\usepackage[utf8]{inputenc}
\usepackage{lmodern}
\usepackage[T1]{fontenc}
\usepackage{microtype}
\usepackage{amsmath,amssymb}
\usepackage{graphicx}
\usepackage{xcolor}
\usepackage{hyperref}
\usepackage{booktabs}
\usepackage{longtable}
\usepackage{listings}
\usepackage{tikz}
\usepackage[alf,abnt-etal-text=it,abnt-etal-list=0]{abntex2cite}

\usetikzlibrary{shapes.geometric, arrows.meta, positioning, shadows, fit}

% --- Configurações de cores ---
\definecolor{blueopencode}{RGB}{37,99,235}
\definecolor{greenopencode}{RGB}{22,163,74}
\definecolor{orangeopencode}{RGB}{234,88,12}
\definecolor{redopencode}{RGB}{220,38,38}
\definecolor{purpleopencode}{RGB}{124,58,237}
\definecolor{graycode}{RGB}{30,41,59}
\definecolor{lightgray}{RGB}{241,245,249}

% --- Configuração de listings ---
\lstdefinestyle{javastyle}{
  language=Java,
  basicstyle=\small\ttfamily,
  keywordstyle=\color{blueopencode}\bfseries,
  commentstyle=\color{green!50!black},
  stringstyle=\color{redopencode},
  numbers=left,
  numberstyle=\tiny,
  frame=single,
  breaklines=true,
  backgroundcolor=\color{lightgray},
  tabsize=2
}

% --- Metadados ABNT ---
\titulo{Pipeline Integrado de Especificação, Validação e Automação no Ecossistema OpenCode: Engenharia de Software Orientada a Especificações com SDD, TDD, Harness e Hooks}
\autor{Equipe OpenCode}
\local{Brasil}
\data{Junho 2026}

\instituicao{%
  OpenCode Ecosystem
  \par
  Engenharia de Software com Agentes Inteligentes
}

\tipotrabalho{Relatório Técnico}

% --- Resumo ---
\begin{document}

\resumo
  Este relatório técnico apresenta a arquitetura integrada de especificação,
  validação e automação do ecossistema OpenCode, fundamentada em quatro pilares:
  \textit{Spec-Driven Development} (SDD), \textit{Test-Driven Development} (TDD),
  infraestrutura de validação (\textit{Harness}) e automação orientada a eventos
  (\textit{Hooks}). A abordagem consolida princípios de Engenharia de Software --
  SOLID, padrões de projeto GoF, \textit{Clean Architecture} e \textit{Continuous
  Delivery} -- em um pipeline fechado de qualidade que garante rastreabilidade
  bidirecional entre especificações e implementação. O ecossistema conta com 186
  componentes documentados, 25 assertivas críticas (CAs) distribuídas em 5
  especificações acadêmicas TDDAcademic, e uma arquitetura de validadores
  orientada a interfaces com suporte a hooks pré, pós e de erro. São discutidas
  as contribuições para as disciplinas de Engenharia de Software, incluindo
  SDD/TDD, CI/CD, arquitetura de software, padrões de projeto e verificação
  formal.
\endresumo

% ======================================================================
\section{INTRODUÇÃO}
% ======================================================================

A engenharia de software moderna demanda abordagens que integrem
especificação rigorosa, validação automatizada e rastreabilidade contínua
entre artefatos \cite{sommerville2011engenharia, pressman2016engenharia}.
O ecossistema OpenCode propõe um pipeline unificado que combina quatro
camadas complementares:

\begin{enumerate}
  \item \textbf{SDD (Spec-Driven Development):} Especificações executáveis
    que funcionam como contratos formais entre requisitos e implementação.
  \item \textbf{TDD (Test-Driven Development):} Assertivas críticas (CAs)
    que validam cada aspecto da especificação.
  \item \textbf{Harness:} Infraestrutura de validação que orquestra a
    execução dos validadores e consolida resultados.
  \item \textbf{Hooks:} Sistema de eventos para automação de ações
    pré e pós-validação.
\end{enumerate}

Este relatório documenta a arquitetura, os padrões de projeto empregados
e as contribuições para as disciplinas de Engenharia de Software,
apresentando o ecossistema OpenCode como estudo de caso de aplicação
sistemática de princípios SDD/TDD em um ambiente multiagente.

% ======================================================================
\section{ARCHITETURA DO PIPELINE}
% ======================================================================

A arquitetura do pipeline segue o fluxo linear com retroalimentação:

\begin{equation}
  \text{SDD} \longrightarrow \text{TDD} \longrightarrow
  \text{Harness} \longrightarrow \text{Hooks}
  \looparrowleft \text{Feedback}
\end{equation}

Cada estágio produz artefatos que alimentam o estágio subsequente,
enquanto o \textit{feedback loop} permite que métricas de cobertura e
resultados de validação retroalimentem a especificação inicial.

% ----------------------------------------------------------------------
\subsection{SDD -- Spec-Driven Development}
% ----------------------------------------------------------------------

O SDD estabelece que toda implementação deve ser precedida por uma
especificação formal. No OpenCode, as especificações são documentos
estruturados que definem:

\begin{itemize}
  \item \textbf{Metadados:} Identificador, categoria, versão, autor.
  \item \textbf{Requisitos funcionais:} Comportamento esperado do sistema.
  \item \textbf{Assertivas Críticas (CAs):} Condições que devem ser
    satisfeitas para validação.
  \item \textbf{Dependências:} Relacionamentos com outras especificações.
\end{itemize}

O ecossistema mantém 186 componentes documentados com 100\% de cobertura
de especificações \cite{OpenCode2025}.

% ----------------------------------------------------------------------
\subsection{TDD -- Test-Driven Development}
% ----------------------------------------------------------------------

O TDD acadêmico no OpenCode materializa-se através do sistema TDDAcademic,
composto por 5 especificações e 25 Assertivas Críticas (CAs):

\begin{longtable}{lll}
  \caption{Especificações TDDAcademic e CAs correspondentes}
  \label{tab:tddacademic}\\
  \toprule
  \textbf{Especificação} & \textbf{ID} & \textbf{CAs} \\
  \midrule
  \endfirsthead
  \toprule
  \textbf{Especificação} & \textbf{ID} & \textbf{CAs} \\
  \midrule
  \endhead
  \bottomrule
  \endfoot
  TDDAcademic Base & SPEC-TDDA-BASE & CA-BASE-001 a 005 \\
  Ciclo de Desenvolvimento & SPEC-TDDA-CICLO & CA-CICLO-001 a 005 \\
  Baby Steps & SPEC-TDDA-BABY & CA-BABY-001 a 005 \\
  Regras de Refatoração & SPEC-TDDA-REGRAS & CA-REGRAS-001 a 005 \\
  Ferramentas TDD & SPEC-TDDA-FERRAMENTAS & CA-FERRAMENTAS-001 a 005 \\
\end{longtable}

Cada CA é classificada por peso (essencial, alta, média) e domain, e
inclui descrição do cenário de validação e procedimento de verificação.

% ----------------------------------------------------------------------
\subsection{Harness -- Infraestrutura de Validação}
% ----------------------------------------------------------------------

A camada Harness é uma infraestrutura de validação composta por quatro
componentes principais que orquestram o ciclo de validação:

\begin{itemize}
  \item \textbf{DiscoveryService:} Varre o classpath para localizar
    validadores registrados via anotação ou convenção de nomenclatura.
  \item \textbf{ExecutionEngine:} Executa cada validador contra a
    especificação alvo, coletando resultados individuais.
  \item \textbf{AggregationService:} Consolida resultados parciais em
    um relatório unificado de validação.
  \item \textbf{ReportingService:} Gera artefatos de saída (JSON, texto,
    HTML) e dispara hooks de pós-validação.
\end{itemize}

O Harness implementa o padrão \textit{Template Method} através da classe
abstrata \texttt{AbstractValidator<T>}, que define o esqueleto do
algoritmo de validação enquanto delega etapas específicas para
subclasses concretas.

% ----------------------------------------------------------------------
\subsection{Hooks -- Automação Orientada a Eventos}
% ----------------------------------------------------------------------

O sistema de hooks implementa o padrão \textit{Observer} \cite{gamma1995design},
permitindo que agentes sejam notificados sobre eventos do pipeline:

\begin{itemize}
  \item \textbf{Pre-Hook (\texttt{on-spec-create}):} Acionado antes da
    criação de uma especificação. Valida a estrutura, verifica
    consistência dos CAs e bloqueia a criação se inválida.
  \item \textbf{Error-Hook (\texttt{on-ca-fail}):} Acionado quando uma
    CA falha. Registra o incidente, notifica agentes responsáveis e
    bloqueia o avanço no pipeline.
  \item \textbf{Post-Hook (\texttt{on-validation-complete}):} Acionado
    após a validação completa. Gera relatório de cobertura, atualiza
    painel de métricas e dispara implantação se aprovado.
\end{itemize}

% ======================================================================
\section{ARQUITETURA OOP E PADRÕES DE PROJETO}
% ======================================================================

A arquitetura de validadores do OpenCode é construída sobre princípios
SOLID \cite{martin2000principles} e padrões de projeto GoF
\cite{gamma1995design}.

% ----------------------------------------------------------------------
\subsection{Hierarquia de Validadores}
% ----------------------------------------------------------------------

O núcleo da arquitetura é definido por uma interface genérica:

\begin{lstlisting}[style=javastyle, caption={Interface IValidator<T>}]
public interface IValidator<T> {
    ValidationResult validate(T target);
    boolean canHandle(Spec spec);
    Category getCategory();
}
\end{lstlisting}

A classe abstrata \texttt{AbstractValidator<T>} implementa o padrão
\textit{Template Method}:

\begin{lstlisting}[style=javastyle, caption={AbstractValidator -- Template Method}]
public abstract class AbstractValidator<T>
    implements IValidator<T> {

    protected String id;
    protected Category category;

    public ValidationResult validate(T target) {
        ValidationResult result = new ValidationResult();
        result.setId(this.id);
        if (!preValidate(target, result)) return result;
        doValidate(target, result);
        postValidate(target, result);
        return result;
    }

    protected boolean preValidate(T t, ValidationResult r) { return true; }
    protected abstract void doValidate(T t, ValidationResult r);
    protected void postValidate(T t, ValidationResult r) {}
}
\end{lstlisting}

As implementações concretas incluem:

\begin{itemize}
  \item \textbf{SpecValidator} (SPEC-VALIDATOR-001): Valida estrutura,
    semântica e dependências de especificações.
  \item \textbf{CAValidator} (CA-VALIDATOR-001): Valida assertivas
    críticas individuais, seus pesos e classificações.
  \item \textbf{HarnessValidator} (HARNESS-VALIDATOR-001): Orquestra a
    execução de múltiplos validadores e gera relatórios de cobertura.
\end{itemize}

% ----------------------------------------------------------------------
\subsection{Padrões de Projeto Empregados}
% ----------------------------------------------------------------------

\subsubsection{Strategy}
Cada validador concreto encapsula um algoritmo de validação específico,
permitindo que sejam intercambiáveis. O \texttt{HarnessValidator}
seleciona dinamicamente o validador adequado com base no tipo da
especificação via \texttt{canHandle()}.

\subsubsection{Template Method}
\texttt{AbstractValidator<T>} define o esqueleto do algoritmo de
validação (\texttt{preValidate → doValidate → postValidate}),
permitindo que subclasses redefinam etapas específicas sem alterar
a estrutura geral.

\subsubsection{Composite}
O \texttt{HarnessValidator} trata validadores individuais e coleções
de validadores de forma uniforme, permitindo composição hierárquica
de validações.

\subsubsection{Observer}
O sistema de hooks (\textit{Publisher}) notifica agentes registrados
(\textit{Subscribers}) sobre eventos do pipeline, implementando
desacoplamento completo entre a origem do evento e os manipuladores.

% ======================================================================
\section{INTEGRAÇÃO DO PIPELINE E CICLO DE QUALIDADE}
% ======================================================================

O pipeline opera como um ciclo fechado de qualidade (\textit{closed-loop
quality}):

\begin{enumerate}
  \item \textbf{Especificação:} Cria-se ou atualiza-se uma especificação
    SDD. O pre-hook valida a estrutura.
  \item \textbf{Implementação:} Desenvolve-se a implementação orientada
    pelas especificações e CAs.
  \item \textbf{Validação:} O Harness executa todos os validadores
    aplicáveis contra a implementação.
  \item \textbf{Feedback:} Os resultados da validação retroalimentam a
    especificação, permitindo refinamento iterativo.
  \item \textbf{Relatório:} O ReportingService gera artefatos de
    cobertura e rastreabilidade.
\end{enumerate}

Este ciclo garante que toda alteração seja validada contra o conjunto
completo de especificações antes de ser considerada concluída,
estabelecendo um contrato formal e verificável entre requisitos e
implementação \cite{humble2010continuous}.

% ======================================================================
\section{EXTRAÇÃO PARA AVALIAÇÃO EM BANCA}
% ======================================================================

Esta seção apresenta os artefatos do OpenCode relevantes para avaliação
em banca acadêmica, mapeados contra disciplinas de Engenharia de Software.

\begin{longtable}{p{4cm}p{4cm}p{4cm}}
  \caption{Mapeamento Artefato-Disciplina}
  \label{tab:banca}\\
  \toprule
  \textbf{Artefato} & \textbf{Disciplinas ES} & \textbf{Contribuição} \\
  \midrule
  \endfirsthead
  \toprule
  \textbf{Artefato} & \textbf{Disciplinas ES} & \textbf{Contribuição} \\
  \midrule
  \endhead
  \bottomrule
  \endfoot
  SDD + 186 specs & SDD, Requisitos, Documentação & Rastreabilidade req.--impl. \\
  TDDAcademic CAs & TDD, Verificação, Testes & 25 assertivas críticas \\
  Harness + Validadores & CI/CD, Arquitetura SWEBOK & Infraestrutura de validação \\
  Hooks + Eventos & Automação, DevOps & Gatilhos de pipeline \\
  Interface IValidator & OOP, SOLID, Padrões GoF & Strategy + Template Method \\
  AbstractValidator & Reuso, Herança, Polimorfismo & Template Method concreto \\
  Diagrama Pipeline & Arquitetura de Software, Documentação & Visão integrada do fluxo \\
  Feedback Loop & Melhoria Contínua, Métricas & Ciclo fechado de qualidade \\
\end{longtable}

% ======================================================================
\section{CONCLUSÃO}
% ======================================================================

O ecossistema OpenCode demonstra a viabilidade de integrar SDD, TDD,
Harness e Hooks em um pipeline unificado de engenharia de software. As
principais contribuições incluem:

\begin{enumerate}
  \item \textbf{Arquitetura orientada a interfaces} com separação clara
    de responsabilidades e adesão aos princípios SOLID.
  \item \textbf{Pipeline de qualidade fechado} que garante rastreabilidade
    bidirecional entre especificações e implementação.
  \item \textbf{Sistema de hooks extensível} que permite automação de
    ações pré e pós-validação sem acoplamento.
  \item \textbf{Documentação completa} com 186 componentes e 100\% de
    cobertura de especificações.
  \item \textbf{Estrutura de validação multiagente} com suporte a
    composição hierárquica de validadores.
\end{enumerate}

Os padrões de projeto empregados -- Strategy, Template Method, Composite
e Observer -- fornecem uma base sólida para extensão futura do
ecossistema, enquanto o ciclo fechado de qualidade estabelece um
arcabouço replicável para projetos que exigem rigor entre especificação
e implementação.

% ======================================================================
\section{TRABALHOS FUTUROS}
% =====================================================================%

Entre as direções para trabalho futuro, destacam-se:

\begin{itemize}
  \item Integração com ferramentas formais de verificação (Z3, Coq) para
    validação simbólica de especificações.
  \item Geração automática de casos de teste a partir das especificações
    SDD utilizando técnicas de \textit{model-based testing}.
  \item Expansão do sistema de hooks para suportar encadeamento de
    eventos e orquestração multiagente.
  \item Métricas quantitativas de eficácia do pipeline em projetos reais.
  \item Publicação do ecossistema como \textit{open-source} para
    replicação e validação pela comunidade acadêmica.
\end{itemize}

% ======================================================================
\bibliography{relatorio_tecnico_opencode}
% ======================================================================

\end{document}
